\section{Landasan Matematis RSA}
\label{sec:landasan-rsa}

Rumus enkripsi dan dekripsi RSA diturunkan dari Teorema Euler. Untuk itu,
fungsi Totient Euler perlu dipahami lebih dahulu.

\subsection{Fungsi Totient Euler}

Fungsi Totient Euler $\varphi(n)$ menentukan berapa banyak bilangan
$1, 2, 3, \dots, n$ yang relatif prima terhadap $n$. Dua bilangan disebut
relatif prima bila pembagi bersama terbesarnya (PBB atau \textit{gcd}) sama
dengan satu.

Sebagai contoh, $\varphi(20) = 8$. Ada delapan bilangan antara 1 dan 20 yang
relatif prima terhadap 20, yaitu
\[
1,\ 3,\ 7,\ 9,\ 11,\ 13,\ 17,\ 19 .
\]

Jika $n = pq$ adalah bilangan komposit dengan $p$ dan $q$ prima, maka
\begin{equation}
\varphi(n) = \varphi(p)\,\varphi(q) = (p-1)(q-1).
\label{eq:totient-rsa}
\end{equation}

\subsection{Teorema Euler}

Teorema Euler menyatakan bahwa bila $a$ relatif prima terhadap $n$, yaitu
$\text{PBB}(a, n) = 1$, maka
\[
a^{\varphi(n)} \equiv 1 \pmod{n}.
\]
Syarat itu perlu diperhatikan. Tanpa syarat tersebut, kekongruenan di atas
tidak dijamin berlaku.

\subsection{Penurunan Rumus Enkripsi dan Dekripsi}

Penurunan rumus RSA bertolak dari Teorema Euler. Kedua ruas dipangkatkan
dengan bilangan bulat $k$, lalu $a$ diganti dengan plainteks $m$:
\begin{align*}
a^{\varphi(n)} &\equiv 1 \pmod{n} \\
a^{k\varphi(n)} &\equiv 1 \pmod{n} \\
m^{k\varphi(n)} &\equiv 1 \pmod{n}.
\end{align*}
Kedua ruas kemudian dikalikan dengan $m$, sehingga
\begin{equation}
m^{k\varphi(n)+1} \equiv m \pmod{n}.
\label{eq:rsa-pangkat}
\end{equation}

Misalkan $e$ dan $d$ dipilih sedemikian sehingga
\begin{equation}
e \cdot d \equiv 1 \pmod{\varphi(n)},
\label{eq:ed-kongruen}
\end{equation}
yang setara dengan
\begin{equation}
e \cdot d = k\,\varphi(n) + 1.
\label{eq:ed-kelipatan}
\end{equation}

Dengan mengganti eksponen pada Persamaan~\eqref{eq:rsa-pangkat} memakai
Persamaan~\eqref{eq:ed-kelipatan}, diperoleh
\begin{equation}
(m^e)^d = m^{e \cdot d} = m^{\,k\varphi(n)+1} \equiv m \pmod{n}.
\label{eq:rsa-bukti}
\end{equation}
Persamaan itu menyatakan bahwa enkripsi lalu dekripsi mengembalikan pesan
semula. Jadi,
\begin{equation}
c = m^e \bmod n, \qquad m = c^d \bmod n.
\label{eq:rsa-rumus}
\end{equation}
